\section*{ВВЕДЕНИЕ}
\addcontentsline{toc}{section}{ВВЕДЕНИЕ}

Актуальность темы исследования обусловлена современными тенденциями в банковском секторе, где эффективность принятия кредитных решений зависит от точности и скорости скоринговых методов. Скоринг является ключевым инструментом для разделения потенциальных заемщиков на «хороших» -- тех, кто с высокой вероятностью вернет кредит, и «плохих» -- тех, кто, скорее всего, не вернет кредит. Эти термины общеприняты в данной области и отражают степень кредитоспособности клиентов \cite{Didikin}.

В скоринге используются косвенные признаки, поскольку прямые признаки надежности заемщика отсутствуют. Оцениваются сведения о заемщике, включая уровень дохода, наличие действующих кредитных обязательств, возраст, трудовой стаж и другие значимые характеристики. На их основе строится комплексный критерий для принятия решения о выдаче кредита.

В последние годы динамика кредитования физических лиц в России характеризуется существенными изменениями. По данным Frank RG, в 2025 году объем выданных кредитов физическим лицам составил около 9{,}89 трлн рублей, что на 25{,}6\% ниже уровня 2024 года. Это стало минимальным значением за последние шесть лет. \cite{FrankRG} Одновременно наблюдается снижение спроса на кредитные продукты на фоне ужесточения денежно-кредитной политики. При этом сохраняется значительный уровень просроченной задолженности, что повышает кредитные риски банковской системы и усиливает требования к точности скоринговых моделей.

Распространение современных финансовых технологий привело к трансформации рынка финансовых услуг и формированию новых моделей финансирования, таких как краудфандинг, краудинвестинг, P2P-кредитование и другие. Для сохранения конкурентоспособности банки вынуждены совершенствоваться, повышая скорость обработки заявок и оценки кредитоспособности заемщиков.


\textbf{Актуальность работы} состоит в необходимости повышения точности и эффективности оценки кредитоспособности заемщиков в условиях роста объемов кредитования и увеличения рисков. Современные банковские системы требуют автоматизированных решений, способных обрабатывать большие объемы данных и учитывать различные факторы, влияющие на вероятность дефолта. Традиционные методы оценки кредитного риска не всегда обеспечивают достаточную точность, что приводит к финансовым потерям. В связи с этим возрастает интерес к методам машинного обучения. Их применение позволяет повысить качество прогнозирования и адаптироваться к изменяющимся условиям.

\textbf{Цель работы} заключается в разработке модели оценки кредитоспособности заемщика на основе методов машинного обучения, обеспечивающей высокую точность прогнозирования кредитного риска.

Для достижения поставленной цели требуется решить следующие задачи:
\begin{itemize}
  \item изучить существующие методы оценки кредитоспособности заемщиков;
  \item провести анализ методов машинного обучения, применяемых в скоринге;
  \item разработать модель оценки кредитного риска на основе градиентного бустинга;
  \item реализовать программную модель и провести ее обучение;
  \item оценить качество разработанной модели и сравнить с альтернативными методами.
\end{itemize}